% LaTeX Template for AI Problem Analysis Output
% AMC12B 2024 Problem 15 - Strategic Analysis
% Framework Version: 2.1 (Enhanced for COMC)

\documentclass[12pt]{article}
\usepackage[utf8]{inputenc}
\usepackage{amsmath, amssymb, amsthm}
\usepackage{geometry}
\usepackage{xcolor}
\usepackage[most]{tcolorbox}
\usepackage{enumitem}
\usepackage{fancyhdr}

% Page setup
\geometry{margin=1in}
\setlength{\headheight}{14.5pt}
\pagestyle{fancy}
\fancyhf{}
\rhead{AMC12 Problem Analysis}
\lhead{2024 AMC 12B Problem 15}
\cfoot{\thepage}

% Color definitions
\definecolor{problemconfig}{RGB}{230, 242, 255}    % Light blue
\definecolor{strategic}{RGB}{255, 230, 242}       % Light pink
\definecolor{tactical}{RGB}{242, 255, 230}        % Light green
\definecolor{techniques}{RGB}{255, 242, 230}      % Light yellow
\definecolor{verification}{RGB}{255, 230, 230}    % Light red
\definecolor{solution}{RGB}{230, 255, 255}        % Light cyan
\definecolor{competition}{RGB}{242, 242, 255}     % Light purple
\definecolor{gold}{RGB}{218, 165, 32}

% Box styles
\newtcolorbox{problemconfigbox}{
    colback=problemconfig,
    colframe=blue,
    title=PROBLEM CONFIGURATION ANALYSIS,
    fonttitle=\bfseries,
    arc=3pt,
    boxrule=1pt,
    breakable
}

\newtcolorbox{strategicbox}{
    colback=strategic,
    colframe=purple,
    title=STRATEGIC FRAMEWORK APPLICATION,
    fonttitle=\bfseries,
    arc=3pt,
    boxrule=1pt,
    breakable
}

\newtcolorbox{tacticalbox}{
    colback=tactical,
    colframe=green,
    title=TACTICAL METHODOLOGY SELECTION,
    fonttitle=\bfseries,
    arc=3pt,
    boxrule=1pt,
    breakable
}

\newtcolorbox{techniquesbox}{
    colback=techniques,
    colframe=orange,
    title=CONVERSION TECHNIQUES APPLICATION,
    fonttitle=\bfseries,
    arc=3pt,
    boxrule=1pt,
    breakable
}

\newtcolorbox{verificationbox}{
    colback=verification,
    colframe=red,
    title=VERIFICATION STRATEGY DESIGN,
    fonttitle=\bfseries,
    arc=3pt,
    boxrule=1pt,
    breakable
}

\newtcolorbox{solutionbox}{
    colback=solution,
    colframe=cyan,
    title=SOLUTION ROADMAP,
    fonttitle=\bfseries,
    arc=3pt,
    boxrule=1pt,
    breakable
}

\newtcolorbox{competitionbox}{
    colback=competition,
    colframe=violet,
    title=COMPETITION STRATEGY INTEGRATION,
    fonttitle=\bfseries,
    arc=3pt,
    boxrule=1pt,
    breakable
}

\newtcolorbox{keyinsightbox}{
    colback=yellow!20,
    colframe=gold,
    title=KEY INSIGHT,
    fonttitle=\bfseries,
    arc=3pt,
    boxrule=2pt,
    leftrule=5pt,
    breakable
}

\newtcolorbox{warningbox}{
    colback=red!10,
    colframe=red!50,
    title=WARNING: COMMON PITFALL,
    fonttitle=\bfseries,
    arc=3pt,
    boxrule=2pt,
    leftrule=5pt,
    breakable
}

\newtcolorbox{tipbox}{
    colback=green!10,
    colframe=green!50,
    title=PRO TIP,
    fonttitle=\bfseries,
    arc=3pt,
    boxrule=2pt,
    leftrule=5pt,
    breakable
}

\begin{document}

\title{AMC12B 2024 Problem 15\\Strategic Problem Analysis}
\author{Competition Math Framework v2.1}
\date{\today}
\maketitle

\section*{Problem Statement}

A triangle in the coordinate plane has vertices $A(\log_2 1, \log_2 2)$, $B(\log_2 3, \log_2 4)$, and $C(\log_2 7, \log_2 8)$. What is the area of $\triangle ABC$?

$$\textbf{(A) }\log_2\frac{\sqrt{3}}{7}\qquad \textbf{(B) }\log_2\frac{3}{\sqrt{7}}\qquad \textbf{(C) }\log_2\frac{7}{\sqrt{3}}\qquad \textbf{(D) }\log_2\frac{11}{\sqrt{7}}\qquad \textbf{(E) }\log_2\frac{11}{\sqrt{3}}$$

\vspace{0.3cm}
\noindent\textit{Note: The answer is not revealed here to allow students to work through the problem. See the Final Answer section at the end.}

\begin{problemconfigbox}
\subsection*{A. Problem Classification}
\begin{itemize}[leftmargin=*]
    \item \textbf{Problem Type}: To Find (specific numerical value)
    \item \textbf{Competition Context}: AMC12B 2024, Problem 15
    \item \textbf{Difficulty Band}: Mid-band (P15) - Late-mid transition, requires multiple technique coordination
    \item \textbf{Mathematical Domain}: Coordinate Geometry with Logarithmic Manipulation
    \item \textbf{Estimated Time}: 3-4 minutes (standard for P15)
    \item \textbf{Point Value}: 6 points (standard AMC12 scoring)
\end{itemize}

\subsection*{B. Problem Structure Identification}
\begin{itemize}[leftmargin=*]
    \item \textbf{Given Information}: 
    \begin{itemize}
        \item Three vertices with logarithmic coordinates: $A(\log_2 1, \log_2 2)$, $B(\log_2 3, \log_2 4)$, $C(\log_2 7, \log_2 8)$
        \item All logarithms have base 2
        \item Coordinates form a triangle in the coordinate plane
    \end{itemize}
    
    \item \textbf{Constraints}: 
    \begin{itemize}
        \item Standard coordinate geometry constraints apply
        \item Answer must be expressed in logarithmic form
        \item Multiple choice format restricts possible answers
    \end{itemize}
    
    \item \textbf{Objective}: Find the exact area of $\triangle ABC$ in logarithmic form
    
    \item \textbf{Hidden Assumptions}: 
    \begin{itemize}
        \item The logarithms can be simplified using basic logarithm properties
        \item The three points actually form a non-degenerate triangle (not collinear)
        \item The area formula should naturally lead to a logarithmic expression
    \end{itemize}
    
    \item \textbf{Answer Choice Leverage}: 
    \begin{itemize}
        \item All answers are in form $\log_2(\text{expression})$
        \item The arguments involve fractions with radicals
        \item The pattern suggests we'll get an expression like $\frac{1}{2}\log_2(\text{something})$
        \item Can use dimensional analysis: area should be positive, so $\log_2(\text{arg})$ could be any sign
    \end{itemize}
    
    \item \textbf{Proof Requirements}: N/A (computational problem, not proof-based)
\end{itemize}

\subsection*{C. Strategic Orientation}
\begin{itemize}[leftmargin=*]
    \item \textbf{Hypothesis}: The logarithmic coordinates can be simplified to standard numbers
    
    \item \textbf{Conclusion}: Need to find area and express it in logarithmic form matching one of the five answer choices
    
    \item \textbf{Notation}: Let the vertices be:
    \begin{align*}
        A &= (x_1, y_1) = (\log_2 1, \log_2 2) \\
        B &= (x_2, y_2) = (\log_2 3, \log_2 4) \\
        C &= (x_3, y_3) = (\log_2 7, \log_2 8)
    \end{align*}
    
    \item \textbf{Intuitive Approach}: 
    \begin{itemize}
        \item First simplify the logarithmic coordinates
        \item Then apply standard area formula for triangle with known coordinates
        \item Finally express result in logarithmic form
    \end{itemize}
    
    \item \textbf{Cross-Domain Potential}: 
    \begin{itemize}
        \item This is primarily coordinate geometry
        \item Requires algebraic manipulation of logarithms
        \item Could potentially use vectors (cross product method)
        \item Shoelace formula is most natural approach
    \end{itemize}
\end{itemize}
\end{problemconfigbox}

\begin{strategicbox}
\subsection*{A. Psychological Strategies}
\begin{itemize}[leftmargin=*]
    \item \textbf{Mental Toughness}: This problem looks intimidating with logarithmic coordinates, but simplifies quickly once you recall $\log_2 2 = 1$, $\log_2 4 = 2$, etc. Don't be deterred by the initial appearance.
    
    \item \textbf{Creativity Cultivation}: Recognize that logarithms of powers of 2 are simple integers. The key insight is that $\log_2(2^n) = n$, so $\log_2 1 = 0$, $\log_2 2 = 1$, $\log_2 4 = 2$, $\log_2 8 = 3$.
    
    \item \textbf{Iterative Refinement}: If the area calculation becomes messy, step back and verify the coordinate simplification. The y-coordinates should be particularly clean (consecutive integers).
\end{itemize}

\subsection*{B. Investigation Strategies}
\begin{itemize}[leftmargin=*]
    \item \textbf{Startup Orientation}: Immediately recognize that $\log_2 1, \log_2 2, \log_2 4, \log_2 8$ are simple values. Simplify coordinates before attempting any area calculation.
    
    \item \textbf{Get Your Hands Dirty}: 
    \begin{itemize}
        \item Write out: $\log_2 1 = 0$, $\log_2 2 = 1$, $\log_2 4 = 2$, $\log_2 8 = 3$
        \item So we have $A(0, 1)$, $B(\log_2 3, 2)$, $C(\log_2 7, 3)$
        \item Notice the y-coordinates are 1, 2, 3 (consecutive integers!)
    \end{itemize}
    
    \item \textbf{Penultimate Step}: The penultimate step is computing the area using the Shoelace formula. The final step is expressing the result in logarithmic form.
    
    \item \textbf{Make It Easier}: Simplify all logarithms that can be simplified immediately. This reduces the problem to a standard coordinate geometry problem.
    
    \item \textbf{Conjecture via Small Cases}: Not directly applicable, but we can verify our final answer makes sense by checking that the area is positive.
    
    \item \textbf{Visualization}: Plot the points approximately:
    \begin{itemize}
        \item $A = (0, 1)$ is on the y-axis
        \item $B \approx (1.58, 2)$ since $\log_2 3 \approx 1.58$
        \item $C \approx (2.81, 3)$ since $\log_2 7 \approx 2.81$
        \item The points form a roughly triangular shape slanting upward to the right
    \end{itemize}
\end{itemize}

\subsection*{C. Late-Band Strategic Playbook}
\begin{itemize}[leftmargin=*]
    \item \textbf{Answer-Choice Leverage}: All answers have the form $\log_2(\text{fraction with radical})$. The area formula will give us something like $\frac{1}{2}|\text{expression}|$, which should convert to $\log_2(\text{something}^{1/2})$ or similar.
    
    \item \textbf{Make It Easier}: The coordinate simplification is the key "make it easier" step. Once we have $A(0,1)$, $B(\log_2 3, 2)$, $C(\log_2 7, 3)$, this becomes a standard problem.
    
    \item \textbf{Pattern Recognition}: The y-coordinates being 1, 2, 3 is a strong hint that the calculation will be clean.
\end{itemize}

\subsection*{D. Argument Construction Strategy}
\begin{itemize}[leftmargin=*]
    \item \textbf{Direct Calculation}: This is a computational problem, so direct calculation is the primary approach.
    
    \item \textbf{Verification Method}: Substitute the answer back or use an alternative area formula (determinant form) to cross-check.
\end{itemize}

\subsection*{E. Cross-Domain Translation}
\begin{itemize}[leftmargin=*]
    \item \textbf{Algebraic to Geometric}: The logarithmic expressions are algebraic, but once simplified, we're working in pure coordinate geometry.
    
    \item \textbf{Coordinate Geometry Visualization}: Drawing a rough sketch helps confirm the triangle is non-degenerate and gives intuition about the expected area magnitude.
\end{itemize}
\end{strategicbox}

\begin{tacticalbox}
\subsection*{A. Core Mathematical Tactics}
\begin{itemize}[leftmargin=*]
    \item \textbf{Primary Tactics}: 
    \begin{itemize}
        \item \textbf{Simplification}: Immediately simplify logarithms of powers of 2
        \item \textbf{Pattern Recognition}: Notice the consecutive integer pattern in y-coordinates
        \item \textbf{Formula Application}: Apply Shoelace theorem or determinant formula for triangle area
    \end{itemize}
    
    \item \textbf{Domain-Specific Tactics}:
    \begin{itemize}
        \item \textbf{Logarithm Properties}:
        \begin{align*}
            \log_2(2^n) &= n \\
            \log_2(ab) &= \log_2 a + \log_2 b \\
            \log_2(a/b) &= \log_2 a - \log_2 b \\
            \log_2(a^k) &= k\log_2 a
        \end{align*}
        
        \item \textbf{Coordinate Geometry}: Shoelace formula for triangle area with vertices $(x_1,y_1)$, $(x_2,y_2)$, $(x_3,y_3)$:
        $$\text{Area} = \frac{1}{2}|x_1(y_2-y_3) + x_2(y_3-y_1) + x_3(y_1-y_2)|$$
    \end{itemize}
    
    \item \textbf{Crossover Tactics}: This problem bridges algebra (logarithms) and geometry (area calculation).
\end{itemize}
\end{tacticalbox}

\begin{techniquesbox}
\subsection*{A. High-Probability Techniques}
\begin{itemize}[leftmargin=*]
    \item \textbf{Primary Technique}: \textbf{Shoelace Formula (M3)} combined with \textbf{Logarithm Identities (M13)}
    
    \item \textbf{Recognition Cues}: 
    \begin{itemize}
        \item Triangle with known coordinates $\Rightarrow$ Shoelace formula
        \item Logarithmic expressions $\Rightarrow$ Simplify using log properties
        \item Answer choices in logarithmic form $\Rightarrow$ Final answer should naturally be a logarithm
    \end{itemize}
    
    \item \textbf{Application}:
    \begin{enumerate}
        \item Simplify all logarithmic coordinates
        \item Apply Shoelace formula: $\frac{1}{2}|x_1(y_2-y_3) + x_2(y_3-y_1) + x_3(y_1-y_2)|$
        \item Express the result in logarithmic form using $\frac{1}{2}\log_2 x = \log_2 \sqrt{x}$
    \end{enumerate}
\end{itemize}

\subsection*{B. Alternative Techniques}
\begin{itemize}[leftmargin=*]
    \item \textbf{Backup Technique 1}: \textbf{Determinant Method}
    \begin{itemize}
        \item Area = $\frac{1}{2}\left|\det\begin{pmatrix} x_1 & y_1 & 1 \\ x_2 & y_2 & 1 \\ x_3 & y_3 & 1 \end{pmatrix}\right|$
        \item This is mathematically equivalent to Shoelace but presented differently
    \end{itemize}
    
    \item \textbf{Backup Technique 2}: \textbf{Vector Cross Product}
    \begin{itemize}
        \item Use vectors $\vec{AB}$ and $\vec{AC}$
        \item Area = $\frac{1}{2}|\vec{AB} \times \vec{AC}|$
        \item More complex for 2D, so Shoelace is preferred
    \end{itemize}
    
    \item \textbf{Backup Technique 3}: \textbf{Geometric Subtraction}
    \begin{itemize}
        \item Create a bounding rectangle
        \item Subtract areas of surrounding triangles/rectangles
        \item This is the approach mentioned in Solution 3 of the problem
    \end{itemize}
    
    \item \textbf{Technique Integration}: The problem requires integrating logarithm simplification with coordinate geometry area formulas.
\end{itemize}
\end{techniquesbox}

\begin{verificationbox}
\subsection*{A. Verification Level Selection}
\begin{itemize}[leftmargin=*]
    \item \textbf{Chosen Level}: Level 4 (Standard Verification) - 30-60 seconds
    
    \item \textbf{Justification}: 
    \begin{itemize}
        \item Problem 15 is mid-band, requiring careful verification
        \item Multiple algebraic steps with logarithms increase error risk
        \item The calculation is straightforward enough that standard verification suffices
        \item We should verify both the coordinate simplification and the area calculation
    \end{itemize}
    
    \item \textbf{Time Budget}: 45 seconds for verification after solving
\end{itemize}

\subsection*{B. Verification Checklist}
\begin{itemize}[leftmargin=*]
    \item \textbf{Coordinate Simplification Check}:
    \begin{itemize}
        \item Verify: $\log_2 1 = 0$ (since $2^0 = 1$) \checkmark
        \item Verify: $\log_2 2 = 1$ (since $2^1 = 2$) \checkmark
        \item Verify: $\log_2 4 = 2$ (since $2^2 = 4$) \checkmark
        \item Verify: $\log_2 8 = 3$ (since $2^3 = 8$) \checkmark
    \end{itemize}
    
    \item \textbf{Computational Accuracy}:
    \begin{itemize}
        \item Check Shoelace formula application
        \item Verify arithmetic in computing $x_1(y_2-y_3) + x_2(y_3-y_1) + x_3(y_1-y_2)$
        \item Check that absolute value is taken correctly
        \item Verify the factor of $\frac{1}{2}$
    \end{itemize}
    
    \item \textbf{Logarithm Conversion}:
    \begin{itemize}
        \item Verify that $\frac{1}{2}\log_2 x = \log_2 \sqrt{x}$ is applied correctly
        \item Check that logarithm properties are used correctly
    \end{itemize}
    
    \item \textbf{Answer Format}:
    \begin{itemize}
        \item Ensure answer matches one of the five given options
        \item Verify the radical is simplified correctly
    \end{itemize}
    
    \item \textbf{Sanity Checks}:
    \begin{itemize}
        \item Area should be positive
        \item The magnitude should be reasonable given the triangle's approximate dimensions
        \item Since $\log_2 3 \approx 1.58$ and $\log_2 7 \approx 2.81$, the area should be approximately $\frac{1}{2} \cdot \text{base} \cdot \text{height} \approx 1$ to $2$ square units in logarithmic space
    \end{itemize}
    
    \item \textbf{Alternative Method Check} (if time permits):
    \begin{itemize}
        \item Recalculate using determinant form
        \item Should get the same result
    \end{itemize}
\end{itemize}

\subsection*{C. Domain-Specific Verification}
\begin{itemize}[leftmargin=*]
    \item \textbf{Geometry Verification}:
    \begin{itemize}
        \item Check that the three points are not collinear (which would give zero area)
        \item The y-coordinates are 1, 2, 3 (distinct), and x-coordinates are 0, $\log_2 3$, $\log_2 7$ (distinct and increasing), so they form a proper triangle
    \end{itemize}
    
    \item \textbf{Algebra Verification}:
    \begin{itemize}
        \item Double-check logarithm arithmetic
        \item Verify that $\log_2(9/7) = \log_2 9 - \log_2 7$
        \item Verify that $\frac{1}{2}\log_2(9/7) = \log_2\sqrt{9/7} = \log_2\frac{3}{\sqrt{7}}$
    \end{itemize}
\end{itemize}
\end{verificationbox}

\begin{keyinsightbox}
\textbf{The Critical Insight}: Immediately simplify logarithms of powers of 2 to integers. This transforms an intimidating logarithmic problem into a straightforward coordinate geometry problem.

Specifically: $\log_2(2^n) = n$, so:
\begin{itemize}
    \item $\log_2 1 = \log_2(2^0) = 0$
    \item $\log_2 2 = \log_2(2^1) = 1$
    \item $\log_2 4 = \log_2(2^2) = 2$
    \item $\log_2 8 = \log_2(2^3) = 3$
\end{itemize}

This reduces the vertices to $A(0, 1)$, $B(\log_2 3, 2)$, $C(\log_2 7, 3)$, where only two coordinates remain in logarithmic form, and the y-coordinates form the clean sequence 1, 2, 3.
\end{keyinsightbox}

\begin{solutionbox}
\subsection*{A. Step-by-Step Solution}

\textbf{Step 1: Initial Setup and Coordinate Simplification}

First, simplify the logarithmic coordinates using the fact that $\log_2(2^n) = n$:

\begin{align*}
    \log_2 1 &= \log_2(2^0) = 0 \\
    \log_2 2 &= \log_2(2^1) = 1 \\
    \log_2 4 &= \log_2(2^2) = 2 \\
    \log_2 8 &= \log_2(2^3) = 3
\end{align*}

The vertices become:
\begin{align*}
    A &= (0, 1) \\
    B &= (\log_2 3, 2) \\
    C &= (\log_2 7, 3)
\end{align*}

\textbf{Checkpoint}: Do these coordinates make sense? Yes - the y-coordinates are consecutive integers (1, 2, 3), and the x-coordinates are increasing (0 $<$ $\log_2 3$ $<$ $\log_2 7$).

\vspace{0.3cm}

\textbf{Step 2: Apply the Shoelace Formula}

The Shoelace formula for the area of a triangle with vertices $(x_1, y_1)$, $(x_2, y_2)$, $(x_3, y_3)$ is:

$$\text{Area} = \frac{1}{2}|x_1(y_2 - y_3) + x_2(y_3 - y_1) + x_3(y_1 - y_2)|$$

Substitute our coordinates:
\begin{align*}
    (x_1, y_1) &= (0, 1) \\
    (x_2, y_2) &= (\log_2 3, 2) \\
    (x_3, y_3) &= (\log_2 7, 3)
\end{align*}

\textbf{Checkpoint}: Is this the right formula? Yes - this is the standard area formula for a triangle given coordinates.

\vspace{0.3cm}

\textbf{Step 3: Calculate Each Term}

Calculate the differences in y-coordinates:
\begin{align*}
    y_2 - y_3 &= 2 - 3 = -1 \\
    y_3 - y_1 &= 3 - 1 = 2 \\
    y_1 - y_2 &= 1 - 2 = -1
\end{align*}

Now compute each term of the sum:
\begin{align*}
    x_1(y_2 - y_3) &= 0 \cdot (-1) = 0 \\
    x_2(y_3 - y_1) &= \log_2 3 \cdot 2 = 2\log_2 3 \\
    x_3(y_1 - y_2) &= \log_2 7 \cdot (-1) = -\log_2 7
\end{align*}

Sum these terms:
$$x_1(y_2 - y_3) + x_2(y_3 - y_1) + x_3(y_1 - y_2) = 0 + 2\log_2 3 - \log_2 7$$

\textbf{Checkpoint}: Does this intermediate result make sense? Yes - we expect a combination of logarithmic terms.

\vspace{0.3cm}

\textbf{Step 4: Simplify Using Logarithm Properties}

Apply logarithm properties:
\begin{align*}
    2\log_2 3 - \log_2 7 &= \log_2(3^2) - \log_2 7 && \text{(Power rule: } k\log_b x = \log_b(x^k)\text{)} \\
    &= \log_2 9 - \log_2 7 \\
    &= \log_2\left(\frac{9}{7}\right) && \text{(Quotient rule: } \log_b a - \log_b c = \log_b(a/c)\text{)}
\end{align*}

\vspace{0.3cm}

\textbf{Step 5: Calculate the Area}

The area is:
\begin{align*}
    \text{Area} &= \frac{1}{2}\left|x_1(y_2-y_3) + x_2(y_3-y_1) + x_3(y_1-y_2)\right| \\
    &= \frac{1}{2}\left|\log_2\left(\frac{9}{7}\right)\right|
\end{align*}

Since $\frac{9}{7} > 1$, we have $\log_2\left(\frac{9}{7}\right) > 0$, so the absolute value equals the expression itself:

$$\text{Area} = \frac{1}{2}\log_2\left(\frac{9}{7}\right)$$

\textbf{Checkpoint}: Is this in the right form? Not quite - the answer choices don't have $\frac{1}{2}$ as a coefficient.

\vspace{0.3cm}

\textbf{Step 6: Convert to the Required Form}

Use the logarithm property $\frac{1}{2}\log_b x = \log_b(x^{1/2}) = \log_b\sqrt{x}$:

\begin{align*}
    \text{Area} &= \frac{1}{2}\log_2\left(\frac{9}{7}\right) \\
    &= \log_2\left(\sqrt{\frac{9}{7}}\right) \\
    &= \log_2\left(\frac{\sqrt{9}}{\sqrt{7}}\right) \\
    &= \log_2\left(\frac{3}{\sqrt{7}}\right)
\end{align*}

\textbf{Final Verification}: This matches answer choice \textbf{(B)}.

\subsection*{B. Verification of Final Answer}

Let's verify our answer using the fact that $\log_2\left(\frac{3}{\sqrt{7}}\right) = \log_2 3 - \frac{1}{2}\log_2 7$:

\begin{align*}
    \log_2\left(\frac{3}{\sqrt{7}}\right) &= \log_2 3 - \log_2(\sqrt{7}) \\
    &= \log_2 3 - \log_2(7^{1/2}) \\
    &= \log_2 3 - \frac{1}{2}\log_2 7
\end{align*}

Multiply by 2:
$$2 \cdot \log_2\left(\frac{3}{\sqrt{7}}\right) = 2\log_2 3 - \log_2 7 = \log_2 9 - \log_2 7 = \log_2\left(\frac{9}{7}\right)$$

This confirms our Shoelace calculation was correct!

\subsection*{C. Alternative Approaches}

\textbf{Alternative Method 1: Determinant Formula}

We can also express the area using the determinant:
$$\text{Area} = \frac{1}{2}\left|\det\begin{pmatrix} 
0 & 1 & 1 \\ 
\log_2 3 & 2 & 1 \\ 
\log_2 7 & 3 & 1 
\end{pmatrix}\right|$$

Expanding along the first row:
\begin{align*}
    \det &= 0 \cdot \begin{vmatrix} 2 & 1 \\ 3 & 1 \end{vmatrix} - 1 \cdot \begin{vmatrix} \log_2 3 & 1 \\ \log_2 7 & 1 \end{vmatrix} + 1 \cdot \begin{vmatrix} \log_2 3 & 2 \\ \log_2 7 & 3 \end{vmatrix} \\
    &= 0 - 1 \cdot (\log_2 3 - \log_2 7) + 1 \cdot (3\log_2 3 - 2\log_2 7) \\
    &= -\log_2 3 + \log_2 7 + 3\log_2 3 - 2\log_2 7 \\
    &= 2\log_2 3 - \log_2 7 \\
    &= \log_2\left(\frac{9}{7}\right)
\end{align*}

Therefore: $\text{Area} = \frac{1}{2}\left|\log_2\left(\frac{9}{7}\right)\right| = \log_2\left(\frac{3}{\sqrt{7}}\right)$ \checkmark

\vspace{0.3cm}

\textbf{Alternative Method 2: Geometric Subtraction}

As mentioned in the problem's Solution 3, we can use a bounding rectangle approach:

\begin{enumerate}
    \item The bounding rectangle has dimensions $\log_2 7$ by $2$ (from $y=1$ to $y=3$).
    \item Its area is $2\log_2 7$.
    
    \item Subtract the areas of surrounding regions:
    \begin{itemize}
        \item Triangle with legs $\log_2 3$ and $1$: Area = $\frac{1}{2}\log_2 3$
        \item Triangle with legs $(\log_2 7 - \log_2 3) = \log_2(7/3)$ and $1$: Area = $\frac{1}{2}\log_2(7/3)$
        \item Rectangle with dimensions $\log_2(7/3)$ and $1$: Area = $\log_2(7/3)$
    \end{itemize}
    
    \item Calculate:
    \begin{align*}
        \text{Area} &= \log_2 7 - \frac{1}{2}\log_2 3 - \frac{1}{2}\log_2(7/3) - \log_2(7/3) \\
        &= \log_2 7 - \frac{1}{2}\log_2 3 - \frac{3}{2}\log_2(7/3)
    \end{align*}
\end{enumerate}

This approach is more complex but arrives at the same answer.
\end{solutionbox}

\begin{competitionbox}
\subsection*{A. Time Management}
\begin{itemize}[leftmargin=*]
    \item \textbf{Estimated Time}: 3-4 minutes total
    \begin{itemize}
        \item Coordinate simplification: 30 seconds
        \item Shoelace formula setup: 30 seconds
        \item Calculation: 1.5-2 minutes
        \item Logarithm conversion: 30 seconds
        \item Verification: 30-45 seconds
    \end{itemize}
    
    \item \textbf{Time Checkpoints}: 
    \begin{itemize}
        \item At 1 minute: Should have simplified coordinates
        \item At 2 minutes: Should be applying Shoelace formula
        \item At 3 minutes: Should have answer in logarithmic form
        \item At 4 minutes: Should be done with verification
    \end{itemize}
    
    \item \textbf{Priority Level}: \textbf{High} - This is a P15 mid-band problem that should be accessible with proper technique. Don't skip it unless completely stuck.
\end{itemize}

\subsection*{B. Risk Assessment}
\begin{itemize}[leftmargin=*]
    \item \textbf{Confidence Level}: 90-95\% after verification
    \begin{itemize}
        \item The coordinate simplification is straightforward
        \item The Shoelace formula is standard and reliable
        \item The logarithm properties are basic
        \item Alternative method confirms the answer
    \end{itemize}
    
    \item \textbf{Common Error Sources}:
    \begin{itemize}
        \item Forgetting to take absolute value in area formula
        \item Sign errors in Shoelace formula
        \item Incorrectly applying logarithm properties
        \item Arithmetic errors in combining logarithmic terms
    \end{itemize}
    
    \item \textbf{Abandonment Criteria}: 
    \begin{itemize}
        \item If stuck after 2 minutes on coordinate simplification $\Rightarrow$ move on
        \item If Shoelace calculation becomes very messy $\Rightarrow$ try alternative method or move on
        \item If can't convert to answer choice form after 4 minutes $\Rightarrow$ guess and flag for review
    \end{itemize}
    
    \item \textbf{Flag for Review}: Only if uncertainty remains after verification, or if time was very tight
\end{itemize}

\subsection*{C. Strategic Context}
\begin{itemize}[leftmargin=*]
    \item \textbf{Problem Positioning}: P15 is in the mid-band range where problems require multiple techniques but are still very accessible with proper approach.
    
    \item \textbf{Score Impact}: This problem is worth 6 points. Getting it right helps maintain momentum through the mid-band.
    
    \item \textbf{Time Budget Implications}: Spending 3-4 minutes here is appropriate and leaves adequate time for P16-20.
    
    \item \textbf{Technique Practice}: This problem reinforces:
    \begin{itemize}
        \item Logarithm properties (essential for AMC12)
        \item Coordinate geometry area formulas
        \item Integration of algebraic and geometric thinking
    \end{itemize}
\end{itemize}
\end{competitionbox}

\begin{warningbox}
\textbf{Common Pitfalls to Avoid}:

\begin{enumerate}
    \item \textbf{Not simplifying coordinates first}: Trying to apply area formulas with unsimplified logarithmic coordinates makes the calculation unnecessarily complex.
    
    \item \textbf{Sign errors in Shoelace formula}: The formula is $x_1(y_2-y_3) + x_2(y_3-y_1) + x_3(y_1-y_2)$. Make sure to subtract in the correct order.
    
    \item \textbf{Forgetting the absolute value}: The area must be positive, so take the absolute value of the determinant/Shoelace result.
    
    \item \textbf{Forgetting the factor of $\frac{1}{2}$}: The Shoelace formula includes a factor of $\frac{1}{2}$ - don't forget it!
    
    \item \textbf{Incorrect logarithm properties}: Remember:
    \begin{itemize}
        \item $k\log_b x = \log_b(x^k)$, not $\log_b(kx)$
        \item $\log_b a - \log_b c = \log_b(a/c)$, not $\log_b(a-c)$
    \end{itemize}
    
    \item \textbf{Not converting to the answer choice form}: The answer choices don't have $\frac{1}{2}$ as a coefficient, so you must convert $\frac{1}{2}\log_2 x$ to $\log_2\sqrt{x}$.
    
    \item \textbf{Arithmetic errors}: Double-check calculations, especially when combining logarithmic terms.
\end{enumerate}
\end{warningbox}

\begin{tipbox}
\textbf{Competition Tips}:

\begin{enumerate}
    \item \textbf{Immediate Simplification}: As soon as you see $\log_2 1$, $\log_2 2$, $\log_2 4$, $\log_2 8$, immediately write down their values (0, 1, 2, 3). This saves time and mental energy.
    
    \item \textbf{Recognize the Pattern}: The y-coordinates being 1, 2, 3 is a strong signal that the problem will have clean calculations.
    
    \item \textbf{Answer Choice Analysis}: All answer choices are positive logarithms with the numerator being either 3, 7, or 11, and the denominator involving $\sqrt{3}$ or $\sqrt{7}$. This suggests we'll get a ratio involving these numbers.
    
    \item \textbf{Logarithm Property Quick Reference}:
    \begin{itemize}
        \item $\frac{1}{2}\log_b x = \log_b\sqrt{x}$
        \item $2\log_b x = \log_b(x^2)$
        \item This is the key conversion needed for this problem
    \end{itemize}
    
    \item \textbf{Numerical Check}: We can verify reasonableness: $\log_2(3/\sqrt{7}) = \log_2 3 - \frac{1}{2}\log_2 7 \approx 1.585 - 1.404 = 0.181$. This positive value is reasonable for a triangle spanning roughly 3 units horizontally ($\log_2 7 \approx 2.81$) and 2 units vertically.
    
    \item \textbf{Verification Strategy}: The quickest verification is to check that your logarithmic answer, when doubled, gives $\log_2(9/7)$ from the Shoelace calculation.
\end{enumerate}
\end{tipbox}

\section*{Final Answer}

After simplifying the logarithmic coordinates to $A(0,1)$, $B(\log_2 3, 2)$, and $C(\log_2 7, 3)$, we apply the Shoelace formula:

$$\text{Area} = \frac{1}{2}|0(2-3) + \log_2 3(3-1) + \log_2 7(1-2)|$$
$$= \frac{1}{2}|2\log_2 3 - \log_2 7|$$
$$= \frac{1}{2}\log_2\left(\frac{9}{7}\right)$$
$$= \log_2\sqrt{\frac{9}{7}}$$
$$= \log_2\frac{3}{\sqrt{7}}$$

$$\boxed{\textbf{(B) } \log_2\frac{3}{\sqrt{7}}}$$

\section*{Confidence Assessment}
\begin{itemize}[leftmargin=*]
    \item \textbf{Confidence Level}: 95\%
    \begin{itemize}
        \item Coordinate simplification is straightforward and verified
        \item Shoelace formula application is standard and double-checked
        \item Logarithm manipulations follow standard properties
        \item Alternative determinant method confirms the same result
        \item Answer matches one of the given choices
    \end{itemize}
    
    \item \textbf{Verification Applied}: Level 4 (Standard Verification) plus alternative method cross-check
    \begin{itemize}
        \item Verified coordinate simplifications
        \item Checked arithmetic in Shoelace formula
        \item Confirmed logarithm property applications
        \item Cross-verified using determinant formula
    \end{itemize}
    
    \item \textbf{Flag for Review}: No - confidence is very high
    
    \item \textbf{Time Spent}: Approximately 3.5 minutes (within target range)
    \begin{itemize}
        \item 30 seconds: Coordinate simplification
        \item 2 minutes: Shoelace formula application and calculation
        \item 30 seconds: Converting to required form
        \item 30 seconds: Verification
    \end{itemize}
\end{itemize}

\section*{Learning Points}

\textbf{Key Takeaways from This Problem}:

\begin{enumerate}
    \item \textbf{Simplification First}: Always simplify expressions before attempting complex calculations. Recognizing $\log_2(2^n) = n$ immediately made this problem tractable.
    
    \item \textbf{Pattern Recognition}: The consecutive integer y-coordinates (1, 2, 3) signaled that the calculation would be clean.
    
    \item \textbf{Shoelace Formula Mastery}: This is a fundamental tool for coordinate geometry problems. Knowing it cold saves significant time.
    
    \item \textbf{Logarithm Properties}: Being fluent with properties like $k\log_b x = \log_b(x^k)$ and $\log_b a - \log_b c = \log_b(a/c)$ is essential.
    
    \item \textbf{Answer Format Matching}: Don't stop until your answer matches the format of the answer choices. The conversion from $\frac{1}{2}\log_2 x$ to $\log_2\sqrt{x}$ was necessary.
    
    \item \textbf{Multiple Methods for Verification}: Having alternative approaches (determinant formula, geometric subtraction) provides strong confidence in the answer.
\end{enumerate}

\vspace{0.5cm}

\textbf{Related Problems to Practice}:
\begin{itemize}
    \item Other AMC12 problems involving logarithmic expressions
    \item Coordinate geometry area problems with algebraic coordinates
    \item Problems requiring integration of algebra and geometry
\end{itemize}

\end{document}

